\begin{afterword}
\summaryhead

The post begins with the author's term, the logical fallacy of generalization from fictional evidence: a journalist writing about AI recalls the \textsc{Terminator} films as if they were a historical case. The inverse error is to be moved too little by history, which really happened. The author suggests that in the ancestral environment there were no movies and what you saw was true, so films move us too much and history, met as ink on paper, too little.

After the mysterious-answer episode, the author realized that vitalism's rise and fall had happened to real people, and that living through such revolutions would probably have prevented the mistake. So the author resolved to think as if everything in the history books had happened to them, since eyes and history books both carry causal traces of events. The world then seemed older and more fragile. The post ends by asking readers to remember being born in a hunter-gatherer tribe, and how often the world has changed in ways they did not guess, so as to be less shocked by the future.

\responsehead

The post names a real contrast: an invented story seen on a screen can move people more than a true one read on a page. What it does with the contrast is weaker.

The explanation is a guess. That films move us too much because ``In our ancestral environment, there were no movies'' is put as a question, ``Is it any wonder,'' and the post gives no evidence for it.

The argument that reading history is like seeing it is also made by a question: ``Is there so much difference'' between the two, since both are causal chains? The post's own footnote shows one difference, that history books over-represent rulers. The paragraph does not take it up, and that footnote is the only caveat the post gives about its method.

The remedy repeats what the diagnosis blamed. The post opened by saying that vividly imagined cases, such as a film, come to be treated as evidence. Its cure for being moved too little by history is to imagine history vividly in the first person: to ``remember'' being born in a tribe, believing you could fly by eating mushrooms, thinking slavery right. The outline is true, but the particulars are the reader's own inventions, and the post does not say how to keep them from counting as evidence in the way the film did.

The motivating claim, that living through past revolutions would ``probably'' have prevented the author's mysterious answer, is hedged. The book's own vitalist, Lord Kelvin, lived through a scientific revolution in his own field and stayed a vitalist (see ``Failing to Learn from History'').

In short: a real contrast between fiction and history, explained by a guess, and answered with an exercise in vivid imagining much like the process the post's opening blamed for the fictional-evidence error.
\end{afterword}
